\begin{abstract}
\large \textbf{Abstract.}
\normalsize
Reinforcement learning (RL) naturally enhances tool-using language agents through execution feedback. However, end-to-end agentic RL is prohibitively expensive due to lengthy tool-use horizons, sequential rollouts, and sparse delayed rewards. In practice, the community relies on supervised fine-tuning (SFT) on expert trajectories, as that is stable and cheap, but suffers from a lack of generalization and lacking explicit reward-driven optimization.

To address these issues, we introduce PivotRL, a turn-level RL framework that bridges SFT and RL using  outcome-verified trajectories that are collected offline. PivotRL formulates agent interactions as an episodic MDP over assistant turns, aligning state transitions with critical tool-call boundaries. Empirically, we find that only a small subset of turns critically influence downstream success. We designate these as pivot states and focus optimization exclusively on them. Crucially, we perform extensive exploration and verification upfront during data collection, thoroughly validating expert actions and exhaustively enumerating all possible pivot states and transition outcomes. Subsequently, we then rely on online RL to solely optimize adherence to these rigorously vetted successful trajectories.

At pivot turns, PivotRL performs a KL-regularized, reward-weighted update toward the outcome-validated expert action, converting long-horizon RL into localized imitation at high-leverage decision points. Uniform rewards recover standard SFT, while non-uniform rewards provide principled policy improvement.

By relying exclusively on logged, extensively verified data, PivotRL enables RL-based improvements from widely available SFT datasets and inference-time interaction traces, removing the need for complex environments or costly rollouts. Through this apporach, PivotRL improves Qwen3-30B-A3B-2507-Thinking from 44.35\% to 62.16\% success rate on Tau-2 (Average of Airline, Retail, and Telecom benchmarks), improves performance from.

\end{abstract}

% \begin{abstract}
%     Reinforcement learning (RL) is the most direct way to improve tool-using language models from execution feedback, but \emph{end-to-end token-level} RL is often impractical: tool-use episodes are long, rollouts are sequential, and standard objectives propagate a single return signal through thousands of tokens even though the environment only reacts at tool/turn boundaries.
% We mitigate this long-horizon bottleneck with two complementary changes.

% First, we formulate agentic tool use as an \emph{episodic Markov decision process at the interaction-turn granularity}, where each step is a complete assistant turn (typically a structured tool call or a short natural-language act) and advantages are defined over turns rather than **only** tokens.
% This allows us to align credit assignment with the interface that actually induces transitions and enabling reuse of cached interaction prefixes instead of repeatedly regenerating full episodes.

% Second, under this turn-level MDP we observe that advantages are typically \emph{sparse}: only a small subset of states exhibit large advantage gaps, while most states have nearly flat action values.
% This motivates \textbf{PivotRL}, which concentrates policy optimization on these high-leverage \emph{pivot states} and applies GRPO-style updates only to the pivot-turn span.
% By combining turn-level credit assignment with pivot-focused updates, PivotRL avoids the need to repeatedly execute full end-to-end sequential rollouts for every gradient step---a key practical barrier for agentic RL in the open source community, rivaling larger models like GPT-OSS-120B.

% On the $\tau^2$ Airline/Retail/Telecom benchmarks, PivotRL improves Qwen3-30B-A3B-2507-Thinking from $44.35$ to $62.16$ average success within $300$ RL steps. Further, the same approach is able to improve Qwen3-30B-A3B-2507-Thinking from $8.75$ success rate to $17.50$ success rate on TerminalBench-Core, using the Terminus harness.
% \end{abstract}
